%%=============================================================================
%% Automatiseren van migratieproces
%%=============================================================================

\chapter{\IfLanguageName{dutch}{Automatisatie}{Automation}}%
\label{ch:automatisatie}

In dit hoofdstuk gaan we kort in op de mogelijkheden om het 
migratieproces
van CA Gen naar PL/I te automatiseren. Hoewel het volledige proces niet
geautomatiseerd kan worden, zijn er verschillende stappen die 
geautomatiseerd kunnen worden om de efficiëntie te verhogen en 
menselijke fouten te verminderen.

\section{Wat kan geautomatiseerd worden en wat niet?}
\label{sec:automatisatie-mogelijkheden}

Waarom kan momenteel niet he volledige migratieproces geautomatiseerd
worden? De belangrijkste reden is dat CA Gen achterliggend heel veel
doet waar de gebruiker geen zicht op heeft. De toolset is als het 
ware een black box. Hierdoor missen automatisatie-tools en AI te veel
context om een volledige migratie te kunnen uitvoeren. Dit zijn een 
paar van de dingen die CA Gen achterliggend doet:
\begin{itemize}
    \item IMS database interacties
    \item SQL interacties
    \item vertalen van bussiness-namen naar technische namen die de 
    mainframe kent
    \item optimalisaties die de toolset zelf doet van flows
\end{itemize}

\subsection{Voorbeeld code generatie door Copilot}

Wat wel kan geautomatiseerd kan worden met AI is het vertalen van de 
pseudo-code in de Action Blocks naar PL/I code. In het onderstaande
voorbeeld hebben we een stukje pseudo-code uit een Action Block genomen
en gevraagd aan Copilot om dit te vertalen naar PL/I code, met als 
context de hoofdstukken van deze bachelorproef die gaan rond de syntax
van PL/I en de mapping van CA Gen constructies naar PL/I constructies.

\begin{center}
    \captionof{listing}{CA Gen pseudo-code die werd vertaald naar PL/I code door Copilot.}
    \label{lst:CA Gen input}
\end{center}

\begin{minted}[fontsize=\small]
EXIT STATE IS sid_verwerking_ab_gestart_rb WITH ROLLBACK

READ mtstbo
    WHERE SOME mat IS CURRENT imp ro
    AND DESIRED mtstbo cod_frm_gg_eh_mat IS EQUAL TO THAT mat cod_frm_gg
    AND DESIRED mtstbo nr_eh_mat_dl IS EQUAL TO THAT mat cod_nr_eh_mat
    AND DESIRED mtstbo nr_eh_mat_dl IS EQUAL TO THAT mat nr_eh_mat
    AND DESIRED mtstbo cod_prc_cre_dl_mat IS EQUAL TO CURRENT imp ro 
    cod_prc_l_ro
    AND SOME matwr is_laatste_proces_voor CURRENT imp ro
    AND THAT matwr nr_wr_eh_mat IS EQUAL TO DESIRED mtstbo 
    nr_prc_cre_dl_mat
    AND DESIRED mtstbo ty_dl_mat IS EQUAL TO "M"

WHEN successful
    WHILE mtstbo cod_eh_mat_dl IS NOT EQUAL TO "W"
        SET loc_imp_mat056 wz54601 cod_frm_gg TO "S"
        SET loc_imp_mat056 wz54601 cod_nr_ro TO "K"
        SET loc_imp_mat056 wz54601 nr_ro TO mtstbo nr_eh_mat_dl
        SET loc_imp_mat056 wz54601 nr_prc_ro TO mtstbo nr_prc_cre_dl_mat
        SET loc_imp_mat056 wz54601 dt_cre_ro TO mtstbo dt_cre_eh_mat_dl
        USE eab_oproep_mat056_lees_prcskin
            WHICH IMPORTS: Work View loc_imp_mat056 wz54601 TO Work View
             imp wz54601
            WHICH EXPORTS: Work View loc_exp_mat056 wz54601 FROM Work 
            View exp wz54601
                             Work View loc_exp_ind_mat056 
                             mat_ind_wz54601 FROM Work View 
                             mat_ind_wz54601
                             Work View loc_exp_wz30911 FROM Work View 
                             exp wz30911

        IF loc_exp_wz30911 cod_tg IS EQUAL TO 1
            AND loc_exp_wz30911 cod_rd IS EQUAL TO 0

* * * * SET gloc_mat sid_werkveld text80 TO concat("Oproep 
eab_oproep_mat055_lees_prcskin gelukt Return: "
* * * * SET gloc_mat_ft int_mat_materiaaleenheid code_frma TO 
loc_imp_mat056 wz54601 cod_frm_gg
* * * * SET gloc_mat_ft int_mat_materiaaleenheid 
codering_identificatienummer TO loc_imp_mat056 wz54601 cod_
* * * * SET gloc_mat_ft int_mat_materiaaleenheid nummer TO 
loc_imp_mat056 wz54601 nr_ro
* * * * USE eab_druk_foutlijn_materiaal
* * * *     WHICH IMPORTS: Group View grl_foutlijn_mat TO Group 
View gri_foutlijn_mat

            IF exp_stamboom_gg int_bsn_bstuyprd pct_lg_skp_na IS 
            EQUAL TO 0
                IF loc_exp_ind_mat056 mat_ind_wz54601 i_lg_skp IS 
                EQUAL TO -1
                ELSE
                    SET exp_stamboom_gg int_bsn_bstuyprd pct_lg_skp_na 
                    TO loc_exp_mat056 wz54601 lg_skp

            IF exp_stamboom_gg int_bsn_bstuyprd brk_emul_skp IS EQUAL 
            TO SPACES
                SET exp_stamboom_gg int_bsn_bstuyprd brk_emul_skp TO 
                loc_exp_mat056 wz54601 emul_brk_skp

        IF mtstbo nr_dl_moe IS NULL
            ESCAPE

        READ mtstbo
            WHERE DESIRED mtstbo cod_dl_mat IS EQUAL TO CURRENT 
            mtstbo cod_dl_moe
            AND DESIRED mtstbo nr_dl_mat IS EQUAL TO CURRENT 
            mtstbo nr_dl_moe

        WHEN successful
        WHEN not found
            ESCAPE

WHEN not found
    EXIT STATE IS mtstbo_nf WITH ROLLBACK

ESCAPE
\end{minted}

\begin{center}
    \captionof{listing}{PL/I output van Copilot op basis van de CA Gen pseudo-code}
    \label{lst:Copilot output}
\end{center}

\begin{minted}[fontsize=\small]
/* ------------------------------------------------------------------ */
/* EXIT STATE IS sid_verwerking_ab_gestart_rb WITH ROLLBACK           */
/* Initialiseer exitstate; WITH ROLLBACK = caller voert ROLLBACK uit  */
/* als dit veld bij terugkeer niet overschreven is.                   */
/* ------------------------------------------------------------------ */

EXP_EXITSTATE = 'SID_VERWERKING_AB_GESTART_RB';

/* ------------------------------------------------------------------ */
/* READ mtstbo                                                        */
/*   WHERE SOME mat is CURRENT imp ro                                 */
/*     AND DESIRED mtstbo ... IS EQUAL TO THAT mat ...                */
/*     AND SOME matur is_laatste_proces_voor CURRENT imp ro           */
/*     AND THAT matur nr_wr_eh_mat IS EQUAL TO DESIRED mtstbo         */
/*              nr_prc_cre_dl_mat                                     */
/*     AND DESIRED mtstbo ty_dl_mat IS EQUAL TO "M"                   */
/*                                                                    */
/* SOME mat / SOME matur = impliciete JOINs via het datamodel;        */
/* DESIRED = veld van mtstbo, THAT = veld van de gejoinde entiteit,   */
/* CURRENT = huidig importview-veld.                                  */
/* ------------------------------------------------------------------ */

EXEC SQL
    SELECT MTSTBO.COD_NR_EH_MAT_DL,
           MTSTBO.NR_EH_MAT_DL,
           MTSTBO.NR_PRC_CRE_DL_MAT,
           MTSTBO.DT_CRE_EH_MAT_DL,
           MTSTBO.COD_DL_MOE,
           MTSTBO.NR_DL_MOE
    INTO  :MTSTBO.COD_NR_EH_MAT_DL,
          :MTSTBO.NR_EH_MAT_DL,
          :MTSTBO.NR_PRC_CRE_DL_MAT,
          :MTSTBO.DT_CRE_EH_MAT_DL,
          :MTSTBO.COD_DL_MOE,
          :MTSTBO.NR_DL_MOE :IND_NR_DL_MOE   /* NULL-indicator */
    FROM  MTSTBO
    INNER JOIN MAT
           ON  MTSTBO.COD_FRM_GG_EH_MAT = MAT.COD_FRM_GG
           AND MTSTBO.COD_NR_EH_MAT_DL  = MAT.COD_NR_EH_MAT
           AND MTSTBO.NR_EH_MAT_DL      = MAT.NR_EH_MAT
    INNER JOIN MATUR
           ON  MATUR.NR_WR_EH_MAT       = MTSTBO.NR_PRC_CRE_DL_MAT
           /* AND <relatie is_laatste_proces_voor: verifieer in 
           datamodel> */
    WHERE  MAT.COD_FRM_GG              = :IMP_RO.COD_FRM_GG
      AND  MAT.COD_NR_EH_MAT           = :IMP_RO.COD_NR_EH_MAT
      AND  MAT.NR_EH_MAT               = :IMP_RO.NR_EH_MAT
      AND  MTSTBO.COD_PRC_CRE_DL_MAT  = :IMP_RO.COD_PRC_L_RO
      AND  MTSTBO.TY_DL_MAT            = 'M';

/* WHEN successful */
IF SQLCODE = 0 THEN DO;

    /* WHILE mtstbo cod_nr_eh_mat_dl IS NOT EQUAL TO "U" */
    /* Label nodig want meerdere ESCAPE-niveaus vanuit de lus */
    LP_MTSTBO:
    DO WHILE (MTSTBO.COD_NR_EH_MAT_DL ^= 'U');
        /* SET loc_imp_mat056 wz54601 cod_frm_gg TO "S" */
        LOC_IMP_MAT056.WZ54601.COD_FRM_GG = 'S';
        /* SET loc_imp_mat056 wz54601 cod_nr_ro TO "K" */
        LOC_IMP_MAT056.WZ54601.COD_NR_RO  = 'K';
        /* SET loc_imp_mat056 wz54601 nr_ro TO mtstbo nr_eh_mat_dl */
        LOC_IMP_MAT056.WZ54601.NR_RO      = MTSTBO.NR_EH_MAT_DL;
        /* SET loc_imp_mat056 wz54601 nr_prc_ro TO mtstbo 
        nr_prc_cre_dl_mat */
        LOC_IMP_MAT056.WZ54601.NR_PRC_RO  = MTSTBO.NR_PRC_CRE_DL_MAT;
        /* SET loc_imp_mat056 wz54601 dt_cre_ro 
        TO mtstbo dt_cre_eh_mat_dl */
        LOC_IMP_MAT056.WZ54601.DT_CRE_RO  = MTSTBO.DT_CRE_EH_MAT_DL;

        /* USE eab_oproep_mat056_lees_prcskin                         */
        /* IMPORTS: loc_imp_mat056 wz54601 → imp wz54601              */
        /* EXPORTS: exp wz54601, mat_ind_wz54601, exp wz30911 terug   */
        CALL EAB_OPROEP_MAT056_LEES_PRCSKIN(
                LOC_IMP_MAT056.WZ54601,
                LOC_EXP_MAT056.WZ54601,
                LOC_EXP_IND_MAT056.MAT_IND_WZ54601,
                LOC_EXP.WZ30911);

        /* IF loc_exp wz30911 cod_tg IS EQUAL TO 1                    */
        /*    AND loc_exp wz30911 cod_rd IS EQUAL TO 0                */
        IF  LOC_EXP.WZ30911.COD_TG = 1
        AND LOC_EXP.WZ30911.COD_RD = 0
        THEN DO;
 
            /* * * * * regels zijn inactief in CA Gen — niet vertalen */

            /* IF exp_stamboom_gg int_bsn_bstyprd pct_lg_skp_na       */
            /*                                IS EQUAL TO 0           */
            IF EXP_STAMBOOM_GG.INT_BSN_BSTYPRD.PCT_LG_SKP_NA = 0 THEN 
            DO;
                /* IF loc_exp_ind_mat056 mat_ind_wz54601 i_lg_skp     */
                /*                                IS EQUAL TO -1      */
                /* → doe niets (lege THEN-tak); ELSE voert de SET uit */
                IF LOC_EXP_IND_MAT056.MAT_IND_WZ54601.I_LG_SKP ^= -1
                THEN
                    /* SET exp_stamboom_gg ... pct_lg_skp_na TO       */
                    /*     loc_exp_mat056 wz54601 lg_skp              */
                    EXP_STAMBOOM_GG.INT_BSN_BSTYPRD.PCT_LG_SKP_NA =
                        LOC_EXP_MAT056.WZ54601.LG_SKP;
            END;

            /* IF exp_stamboom_gg int_bsn_bstyprd brk_emul_skp        */
            /*                          IS EQUAL TO SPACES            */
            IF EXP_STAMBOOM_GG.INT_BSN_BSTYPRD.BRK_EMUL_SKP = ' '
            THEN
                /* SET exp_stamboom_gg ... brk_emul_skp TO            */
                /*     loc_exp_mat056 wz54601 emul_brk_skp            */
                EXP_STAMBOOM_GG.INT_BSN_BSTYPRD.BRK_EMUL_SKP =
                    LOC_EXP_MAT056.WZ54601.EMUL_BRK_SKP;

        END; /* IF cod_tg = 1 AND cod_rd = 0 */

        /* IF mtstbo nr_dl_moe IS NULL → ESCAPE (verlaat de WHILE)    */
        /* IND_NR_DL_MOE is de NULL-indicator van het meest recente 
        READ */
        IF IND_NR_DL_MOE < 0
        THEN LEAVE LP_MTSTBO;

        /* READ mtstbo                                                */
        /*   WHERE DESIRED mtstbo cod_dl_mat IS EQUAL TO              */
        /*             CURRENT mtstbo cod_dl_moe                      */
        /*     AND DESIRED mtstbo nr_dl_mat IS EQUAL TO               */
        /*             CURRENT mtstbo nr_dl_moe                       */
        /* Volg de keten: laad het "moeder"-record.                   */
        EXEC SQL
            SELECT COD_NR_EH_MAT_DL,
                   NR_EH_MAT_DL,
                   NR_PRC_CRE_DL_MAT,
                   DT_CRE_EH_MAT_DL,
                   COD_DL_MOE,
                   NR_DL_MOE
            INTO  :MTSTBO.COD_NR_EH_MAT_DL,
                  :MTSTBO.NR_EH_MAT_DL,
                  :MTSTBO.NR_PRC_CRE_DL_MAT,
                  :MTSTBO.DT_CRE_EH_MAT_DL,
                  :MTSTBO.COD_DL_MOE,
                  :MTSTBO.NR_DL_MOE :IND_NR_DL_MOE
            FROM  MTSTBO
            WHERE COD_DL_MAT = :MTSTBO.COD_DL_MOE
              AND NR_DL_MAT  = :MTSTBO.NR_DL_MOE;

        /* WHEN successful → lus gaat verder                          */
        /* WHEN not found  → ESCAPE (verlaat de WHILE)                */
        IF SQLCODE = 100
        THEN LEAVE LP_MTSTBO;
    END LP_MTSTBO; /* einde WHILE mtstbo cod_nr_eh_mat_dl ^= 'U' */
END; /* WHEN successful (outer READ) */

/* WHEN not found (outer READ)                                        */
ELSE IF SQLCODE = 100 THEN DO;
    /* EXIT STATE IS mtstbo_nf WITH ROLLBACK                          */
    /* Caller controleert EXP_EXITSTATE en voert EXEC SQL ROLLBACK uit*/
    EXP_EXITSTATE = 'MTSTBO_NF';
    RETURN; /* ESCAPE = verlaat dit Action Block */
END;
\end{minted}

\begin{center}
    \captionof{listing}{CA Gen pseudo-code die werd vertaald naar PL/I code door expert}
    \label{lst:expert output}
\end{center}

\begin{minted}[fontsize=\small]
/* ------------------------ */
BEPAAL_SKP_GEGS_VIA_MTSTBO:
PROC;

  #TRACE_ENTRY;

  /* Init variabelen */
  LG_SKP = 0;
  BRK_EMUL_SKP = '';

  H_MTSTBO.COD_FRM_GG_EH_MAT  = H_RO_IMP.COD_FRM_GG_RO;
  H_MTSTBO.COD_NR_EH_MAT_DL   = H_RO_IMP.COD_NR_RO;
  H_MTSTBO.NR_EH_MAT_DL       = H_RO_IMP.NR_RO;
  H_MTSTBO.COD_PRC_CRE_DL_MAT = H_RO_IMP.COD_PRC_L_RO;
  H_MTSTBO.NR_PRC_CRE_DL_MAT  = H_RO_IMP.NR_PRC_L_RO;
  H_MTSTBO.TY_DL_MAT          = 'M';  /* MATDEEL */

  CALL SQL_OPEN_CURSOR('MTSTBO_MATDL_PRC_L');
  CALL SQL_FETCH_CURSOR('MTSTBO_MATDL_PRC_L');

  SCAN_LAATSTE_PROCES:
    DO WHILE(SQLCA.SQLCODE = 0);

    VERWERK_MAT_DEEL:
      DO WHILE(SQLCA.SQLCODE = 0 & H_MTSTBO.COD_NR_EH_MAT_DL = 'K');

        CALL OPROEP_MAT056;

        IF WZ30911.COD_TG = 1 & WZ30911.COD_RD = 0
        THEN
        DO;
          IF LG_SKP = 0
          THEN
            IF WZ54601.GRP_EXP_IND.LG_SKP ^= -1
            THEN
              LG_SKP = WZ54601.GRP_EXP.LG_SKP;

          IF BRK_EMUL_SKP = ''
          THEN
            BRK_EMUL_SKP = WZ54601.GRP_EXP.EMUL_BRK_SKP;
        END;

        IF I_MTSTBO.NR_DL_MOE >= 0
        THEN
          CALL SELECT_MTSTBO_MOEDER;
        ELSE
          LEAVE VERWERK_MAT_DEEL;

    END VERWERK_MAT_DEEL;

    CALL SQL_FETCH_CURSOR('MTSTBO_MATDL_PRC_L');

    END SCAN_LAATSTE_PROCES;

    CALL SQL_CLOSE_CURSOR('MTSTBO_MATDL_PRC_L');
\end{minted}

\section{Verschillen}
\label{sec:verschillen-in-automatisatie}

Er zitten heel wat verschillen tussen de AI-gegenereerde code en de code van de expert. Het grootste verschil dat direct opvalt is de naamgeving van alle variabelen. Dit is ook logisch want developpers programmeren met business namen die CA Gen automatisch aan de hand van de host-encyclopedie omzet naar technische namen die de mainframe wel kent. Ook deelt de expert logica op over verschillende procedures, vooral SQL-blokken, om code leesbaarder te maken.

\section{Conclusie}
\label{sec:automatisatie-conclusie}

Het volledige migratieproces van CA Gen naar PL/I kan momenteel niet
volledig worden geautomatiseerd. Wel kunnen specifieke stappen zoals
pseudo-codevertaling, SQL-mapping en herhalende routines ondersteund
worden met tooling en AI om de efficiëntie te verhogen en fouten te
verminderen.

AI kan een nuttige hulp zijn bij het genereren van PL/I-code en het
verduidelijken van complexe CA Gen-constructies, maar menselijke
controle blijft noodzakelijk voor context, datamodelrelaties en
mainframe-optimalisaties. Daarom blijft de meest realistische aanpak een
hybride proces: automatisering ondersteunt de migratie, terwijl experts
de resultaten verifiëren en bijsturen.

